% I. Предметна област и литературен преглед.
% Обхваща темите от учебната програма: моделиране и споделяне на информация,
% синтаксис и семантика, информационен модел на семантичния уеб, описание на
% ресурси с RDF, обмен на информация с RDF, семантика в уеб пространството,
% онтологии и онтологични системи, OWL, SPARQL, добавяне на правила.
\section{Предметна област и литературен преглед}
\label{sec:literature}

\subsection{Информационен модел на семантичния уеб}

Класическият уеб свързва документи. Връзката между тях е хипервръзка без тип: тя
казва, че два документа са свързани, но не и как. Семантичният уеб добавя точно
този липсващ пласт, като описва не документите, а ресурсите, и дава на всяка връзка
собствен идентификатор и значение. Единицата на записа е триплет от подлог,
сказуемо и допълнение. Множество триплети образува насочен етикетиран граф, в който
допълнението на един триплет може да бъде подлог на друг, и така отделно съставени
описания се свързват без предварително договаряне на обща схема.

Абстрактният синтаксис на този модел е дефиниран в спецификацията RDF~1.1, която
въвежда трите вида термини, тоест IRI, литерал и празен възел, и определя графа и
множеството от графи~\cite{w3c_rdf11_concepts}. Спецификацията умишлено разделя
абстрактния модел от конкретните му записи. Това разделение е важно за настоящия
проект: то позволява синтактичните анализатори да се разглеждат като сменяеми
входове към един и същ вътрешен модел.

\subsection{Синтаксиси за обмен на RDF}

Обменът на RDF между системи става чрез конкретен синтаксис. Turtle е компактен
запис, четим от човек, с префикси, съкратени форми за списъци и вложени описания на
празни възли~\cite{w3c_turtle}. N-Triples е подмножество на Turtle, в което всеки
ред съдържа точно един пълен триплет без съкращения~\cite{w3c_ntriples}. Заради тази
простота N-Triples е удобен за поточна обработка и за тестови набори, докато Turtle
е предпочитан за ръчно писани описания. Проектът поддържа и двата, като N-Triples се
третира като частен случай в същия синтактичен анализатор.

\subsection{Семантика, онтологии и правила}

Синтаксисът определя кои записи са допустими, но не и какво следва от тях.
Семантиката на RDF задава кога един граф следва от друг, тоест отношението на
извеждане~\cite{w3c_rdf11_semantics}. RDF Schema надгражда с речник за класове,
свойства, области и обхвати, и с това позволява прости правила от вида „ако ресурс
принадлежи на подклас, той принадлежи и на надкласа“~\cite{w3c_rdf_schema}. Езикът
OWL~2 разширява изразителната сила значително, като добавя ограничения върху
свойствата, кардиналности, дизюнктни класове и други конструкции, и с това променя
изчислителната сложност на извеждането~\cite{w3c_owl2_overview}. Настоящият проект
реализира извеждане по RDFS като отделен слой и разглежда OWL~2 само описателно.

\subsection{Език за заявки и подходи за индексиране}

SPARQL~1.1 е нормативният език за заявки над RDF. Той дефинира базови графови
шаблони, незадължителни части, обединения, филтри, агрегация, подзаявки и
транслативни пътища по свойства, а също и алгебра, към която текстът на заявката се
превежда преди изпълнение~\cite{w3c_sparql11_query}. Тази алгебра е удобната граница
между синтактичния анализатор и механизма за изпълнение и е приета за такава и в
този проект.

Литературата за съхранение на триплети се съсредоточава върху един въпрос: базовият
графов шаблон изисква многократно търсене по различни комбинации от свързани и
свободни позиции, а един ред на сортиране обслужва добре само част от тях. Hexastore
предлага пълния набор от шест пермутации на подлог, сказуемо и допълнение, с което
всяка комбинация се обслужва от подходящо подреден индекс, срещу цената на шесткратно
повторение на данните~\cite{weiss2008hexastore}. RDF-3X показва, че при речниково
кодиране на термините и компресирани индекси този подход е практически приложим, и
добавя планировчик, който избира реда на съединенията по статистики върху самите
индекси~\cite{neumann2010rdf3x}. Двете публикации формират основата на проектното
решение в Раздел~\ref{sec:design}.

\TODO{преглед на поне два по-нови източника за изпълнение на SPARQL, за да не е
прегледът съсредоточен само в периода 2008 до 2010}
